\chapter{Introduction}
\label{ch:introduction}

% NARRATIVE GOAL OF THIS CHAPTER
% Make the reader care, then promise exactly what the thesis delivers. A bachelor thesis is
% judged partly on whether the question is clear by the end of page three.

\section{Motivation}
\label{sec:motivation}
% Industrial mobile robots move material between machines. The interesting decision is not "how do
% I drive there" (solved by navigation stacks) but "what is worth doing next": a full buffer stops
% a production line, a flat battery strands the robot, and every metre costs energy.
% Keep it concrete: the factory of this thesis produces 20, 30 and 5 units per hour in sections
% A, B and C; a buffer that fills is production lost for good.

\section{Problem statement}
\label{sec:problem-statement}
% One robot, three producing sections, one delivery point, one charger. A decision maker must pick
% the next action under hard constraints (battery reserve, motor temperature, payload, condition).
% Two families of answer exist: hand-written rules (predictable, but blind to trade-offs) and
% learned policies (good at trade-offs, but with no guarantee). This thesis asks whether the two
% can be combined so that the guarantee survives.

\section{Research questions}
\label{sec:research-questions}
% State them as questions you actually answer in \cref{ch:results}. Suggested:
\begin{enumerate}[label=\textbf{RQ\arabic*.}, leftmargin=*]
  \item \textit{<Does a neuro-symbolic policy outperform a deep reinforcement learning policy of
        equal information on throughput, energy and safety?>}
  \item \textit{<What does each half contribute — do the rules alone, or the network alone, suffice?>}
  \item \textit{<Do the policies still hold up in situations they were never trained on?>}
  \item \textit{<Do results obtained in a fast abstract simulator transfer to full physics
        simulation with a real navigation stack?>}
\end{enumerate}

\section{Contributions}
\label{sec:contributions}
% Be precise and modest; every item must be defensible with a number from \cref{ch:results}.
\begin{itemize}[leftmargin=*]
  \item A complete, reproducible ROS~2 system in which a single robot works an autonomous shift:
        factory model, energy/thermal/wear model, mission executor with recovery, decision service
        and a read-only monitor.
  \item A \NS whose hard rules can never be overruled, whose priority tiers are explicit, and whose
        learned ranker is \emph{bounded}: the correction it may apply is limited to
        $\pm\maxCorrection$ score units.
  \item The observation that motivates that bound: without it, the learned ranker carried a habit
        from its training scenarios into an overloaded factory it had never seen and waited while
        production lines stopped (\heavyUnbounded{} against \heavyBounded{} with the bound).
  \item An evaluation of 3240 simulated shifts with paired statistics and a held-out set of
        scenarios, validated by full Gazebo shifts in which the energy model agrees to
        \gazeboEnergyGap{}.
  \item A planner study showing that a global planner selected on short benchmark tours was the
        worse choice over full shifts (\thetaFaults{} shifts faulty against \navfnFaults{}).
\end{itemize}

\section{Outline}
\label{sec:outline}
% One or two sentences per chapter. Write this LAST, when the chapters exist.
